\section{Jurisprudence and Judicial Optimisation}
\label{p13:sec:judicial}

\noindent
The paper turns now from the analysis of trust to its practice, and takes the institutional scale first, the scale at which the problem appeared in the symptomatology as the overload of justice and the sword that hangs but does not fall. If force is not the cause of trust but only a failing device for the production of expectation, then the institutions built upon force\,---\,the courts, the apparatus of enforcement\,---\,stand in need of a reconception, and this section asks what it would be. Its answer has a spine that is counter-intuitive and will be stated plainly: that the effective optimisation of the juridical, in an age when force has failed as a signal, is the \emph{decentralisation of enforcement}\,---\,not the strengthening of the sword but the distribution of the production of trust away from the single channel of compulsion. But the section must guard, from the outset, against the way this spine could be misheard, and the guard determines the whole method of the chapter.

For the spine could be misheard as an idealist exhortation\,---\,as a call to trust more, to cultivate reputation, to repair rather than to punish, as though the institutional form of trust were a matter of enlightened choice, a better idea that societies might adopt if only they saw its merits. The series is historical-materialist, and must treat the matter dialectically; and an idealist exhortation is precisely what the historical-materialist method forbids, for it would mistake the superstructure for something freely adjustable apart from its material base. The institutional forms of trust are not chosen; they are thrown up by the material conditions of their production, and they change as those conditions change. The chapter therefore proceeds not by recommending a reform but by analysing a transformation already under way\,---\,its material causes, its emergent tendencies, and the new contradiction into which it is, under present relations of production, being driven.

\subsection{The double bind of the juridical}
\label{p13:sub:jud-doublebind}

The condition to be analysed is, first, a double bind, in which the two faces of the symptomatology\,---\,the overload and the failure of force\,---\,compose a vicious circle. As the relational and communal orders of trust thin, more of the burden of holding people to their dealings falls upon the courts (the overload); the courts, asked to produce by compulsion an order they can only ever protect, find their signal of compulsion diluted across a volume it cannot cover, and force fails as a signal (the failure); and as force fails, trust thins further, throwing yet more burden upon the courts. The more the juridical is relied upon to generate trust, the more its single signal is diluted; the more diluted the signal, the less it generates; the less it generates, the more it is relied upon. The juridical is caught being asked to \emph{produce}, by force, the very thing\,---\,trust\,---\,that \xref{p13:sec:mechanism} showed force cannot produce but only, at most, signal; and asked to produce it at a scale at which even the signalling fails. This is the double bind, and no increase in the force applied can resolve it, for the problem is not a shortage of force but the exhaustion of force as a means to this end.

\subsection{The counter-intuitive thesis: decentralisation}
\label{p13:sub:jud-decentralisation}

The resolution is not to add force but to cease asking the single channel of force to do what only the coupling of several channels can do. \xref{p13:sec:mechanism} established that well-founded trust is produced by the multiplicative coupling of channels, and that its robustness lies not in the strength of any one but in the improbability of all being forged together. The juridical optimisation that follows from this is the institutional application of the same principle: distribute the production of trust away from the single channel of compulsion and across the several channels the mechanism identified, so that trust is generated by their coupling rather than by the sword alone. This is what decentralisation means here, and why it is effective where the strengthening of force is not. A single channel can be evaded, diluted, and\,---\,under the crisis of the symbol\,---\,collapsed in its precision-weight across the board; several coupled channels, whose mutual consistency is hard to forge or evade together, are robust exactly where the single channel is fragile. The thesis is counter-intuitive because the reflex, faced with a failure of enforcement, is to call for more enforcement; the mechanism shows this to be the marginally decreasing application of a means already exhausted, and the dispersal of trust-production across coupled channels to be the structurally sounder course.

But\,---\,and here the historical-materialist method asserts itself against any idealist reading\,---\,this decentralisation is not a reform to be recommended; it is a transformation already under way, forced by a change in the material conditions of the production of trust, and the rest of the chapter must show it as such.

\subsection{The transformation, dialectically}
\label{p13:sub:jud-dialectic}

The transformation must be grasped in four moments, and the fourth forbids the celebration the first three might invite.

\subsubsection{The material basis of centralised force}
\label{p13:sub:jud-thesis}

Centralised, force-based justice\,---\,the sword wielded by the state\,---\,was not a bad idea that societies happened to adopt; it was the form of trust-production appropriate to, and thrown up by, a definite material base. In a society of strangers, in which exchange is predominantly market exchange and persons are atomised and rendered mobile and interchangeable by the commodity form, trust can no longer be produced by the face-to-face relational history and communal witness that produced it in smaller and more settled worlds. Where the village and the lineage no longer hold a person to his word, the holding must be outsourced to an abstract, impersonal third party\,---\,the state, the law, the sword\,---\,which can compel performance among strangers who have no other hold upon one another. Centralised force was, that is, the historically necessary form of trust-production adequate to the material base of a market society of strangers. This must be said first, on pain of turning the critique into a moralising lament: the sword was not an error but an adaptation, and a correct one, to its conditions.

\subsubsection{The material failure}
\label{p13:sub:jud-antithesis}

Its failure, likewise, is material and not a decay of values. The double bind of \xref{p13:sub:jud-doublebind} is not produced by a hardening of hearts or a decline of civic virtue; it is produced by a change in the material conditions. The scale and speed of exchange, under globalisation and digitalisation, have outrun the physical capacity of any apparatus of enforcement; the volume of breach has outrun the carrying capacity of the courts; and the cost of producing a convincing symbol has fallen to zero, saturating the environment with forged signals and collapsing, as \xref{p13:sub:mech-crisis} showed, the precision-weight a population rationally assigns to the symbolic channel on which centralised enforcement partly relied. These are changes at the level of the forces and the mode of production, and they have brought the once-adequate form of trust-production\,---\,centralised force\,---\,into contradiction with the new material conditions. This is the historical-materialist account of the crisis of the symbol: not the idealist story of a value-nihilism that has overtaken the culture, but the materialist story of a contradiction between the established mode of producing trust and the changed material base beneath it. The diagnosis of \xref{p13:sec:diagnosis} is here raised from a cultural diagnosis to a historical-materialist one.

\subsubsection{The emergent recomposition}
\label{p13:sub:jud-emergence}

The dialectical point follows, and it is the one that distinguishes analysis from exhortation. The decentralisation of enforcement is not something we ought to choose; it is a recomposition that the contradiction is already forcing into being. The new forms of non-coercive trust-production\,---\,reputation systems, credit mechanisms, platform ratings, restorative justice, the distributed ledgers of trust\,---\,are not the good ideas of reformers; they are the spontaneous symptoms of a mode of trust-production reorganising itself under new material conditions, the new technologies of network and computation throwing up new channels as the old single channel fails. The task of the section is therefore not to \emph{recommend} decentralisation but to \emph{analyse} this recomposition that is already occurring: to identify its direction, its potential, and\,---\,the fourth moment\,---\,its internal contradiction.

\subsubsection{The capture: the new contradiction}
\label{p13:sub:jud-capture}

For the dialectical method requires that the emergent form be examined for its own contradictions, and not celebrated. The new channels of non-coercive trust are, under the present relations of production, being captured by capital and turned to alienation. Reputation and credit, which might be the emancipatory institutionalisation of the relational-historical channel, become under capital the social-credit score, the algorithmic management of the worker, credit as a new instrument of discipline\,---\,a finer and more pervasive control than the sword, precisely the retreat to the quantifiable that \xref{p13:sub:sympt-universal} described, with the metrics now held by the platform and the owner. The platform's badge of trust becomes capital's commodification of trust, a thing the platform can monopolise and extract rent upon. So the decentralised technical form, under centralised ownership of the new channels, may produce a domination more insidious than the sword it replaces\,---\,a domination that operates through the very channels that were to have freed trust from compulsion. The question the historical-materialist method therefore presses, and that the justice criterion of \xref{p13:sec:justice} presses with it, is: who owns these new channels, and for whose benefit do they run\,---\,to whom does the value they produce return, and is there one who is rendered, in them, mere fuel? The decentralisation of the \emph{technical form} of trust-production is no emancipation in itself; under the wrong relations of production it is a new mode of extraction. The synthesis is therefore not the triumphant slogan that decentralisation has freed us, but a new contradiction: an emergent form whose emancipatory potential is real and whose capture by capital is equally real, and whose outcome is not settled by the technology but by the relations of production that own it.

\subsection{The three channels institutionalised}
\label{p13:sub:jud-channels}

Within that dialectical frame, and not as free-standing recommendations, the three non-coercive channels of \xref{p13:sec:mechanism} can be seen taking institutional form, each with its emancipatory possibility and each with its characteristic danger of capture.

Reputation and credit systems are the institutionalisation of the relational-historical channel: they aggregate the scattered relational histories of \xref{p13:sec:relational} into a transmissible, institutional form of inductive trust, making a person's accumulated track of dealing legible beyond the circle that witnessed it. This is their power\,---\,the institutional aggregation of inductive trust\,---\,and their danger is exactly the one named above: the degeneration of the credit record into a purely quantitative instrument of discipline, the relational history flattened into a score and the score turned to control.

Restorative justice is the institutionalisation of the self-committing channel. Rather than imposing a sanction by force, it creates the occasion for an offender's real, irrevocable, self-implicating assumption of what he has done\,---\,the encounter with the one he has wronged, the genuine and costly taking-on of responsibility\,---\,which produces, through the high-fidelity likelihood of a real self-binding rather than through compulsion, the rebuilding of an expectation that a mere unenforceable judgment could not~\parencite{zehr2002}. It works the self-committing channel of \xref{p13:sub:mech-likelihood} at the scale of the institution; and its danger is the hollowing of the encounter into a procedural ritual, a staged contrition that forges the self-commitment it should embody.

Public registration and transparency are the institutionalisation of the projective path of \xref{p13:sub:trans-projection}. They do not manufacture trust; they let a real state of performance\,---\,a real record of what has and has not been done\,---\,cast, in the symbolic register the evidentialist reads, the shadow by which it can be recognised. They make the real visible rather than producing the appearance of it; and their danger is the surveillance that the same transparency enables, the registry turned from a means of letting the real be seen into a means of total observation.

\subsection{The place of the sword: relocation, not abolition}
\label{p13:sub:jud-sword}

A boundary must be drawn as firmly here as the spine of the chapter is asserted, for the spine is liable to a dangerous over-reading, and the historical-materialist honesty that refuses to celebrate the emergent form refuses equally to romanticise the abandonment of force. The decentralisation of trust-production is \emph{not} the abolition of enforcement. Force retains an indispensable role\,---\,but it is a different role from the one it was failing at. Force is not, and never was, the producer of trust; but it is the necessary backstop, the after-the-fact protection, and above all the protection of the vulnerable against predation and against the betrayal that the powerful might otherwise inflict with impunity. The weak still need the sword\,---\,not to generate the trust on which they deal, which force cannot generate, but to protect them when trust is betrayed and to guarantee that the betrayal of the powerful is not costless. Decentralisation of trust-production is not de-coercion; the sword is not taken away but \emph{relocated}, from the false position of the producer of trust, which it cannot occupy, to its true position as the backstop that protects against the worst, especially on behalf of those least able to protect themselves. To forget this would be to abandon, in the name of a purer trust, exactly those whom the purer trust does not yet protect.

\subsection{The juridical as the practice of a relational nomos}
\label{p13:sub:jud-nomos}

The chapter closes by connecting its analysis to the foundational paper and to the series, for the transformation it has traced is the institutional form of a movement the foundational paper named at the largest scale. That paper, drawing on the line of Marx and of the theorists of social reproduction, marked the overload of justice as a diagnosis without an exit, and held open the thought\,---\,against the horizon of the eventual withering of the state\,---\,that the law need not wither with it, but might pass from a law of action, the law of the sword that compels, toward a law of non-action, a law of pure relation, a relational \emph{nomos}~\parencite{wanhong_grb,marx1976}. The present chapter supplies the exit the foundational paper marked but did not develop. The centralised, force-based law\,---\,the law of action, the sword\,---\,decays as its material base dissolves; and a nomos built upon the production of relational trust\,---\,the law of non-action, the law of pure relation\,---\,germinates in the contradiction. But whether it is realised, or captured, depends, as the whole chapter has insisted, not upon the technical form but upon whether the relations of production are themselves transformed\,---\,upon who comes to own the new channels of trust, and upon whether the value they generate is returned, justly, to those whose dealing generates it, with no one rendered fuel. The optimisation of the juridical is, in the end, the freeing of justice from the impossible burden of being the last Other, the final guarantor of an order it was never able to underwrite by force; and the return of justice to its place as the practice, and the protector, of a relational nomos whose trust is generated elsewhere\,---\,in the coupled channels, at the scale of relation, that the rest of the paper has described. To that scale, the scale of a life and its nearest bonds, the paper now turns.
